% ---------------------------------------------------------------------------
% Capítulo de requisitos e backlog do protótipo.
% Fontes: docs/requisitos-app.md, docs/backlog/backlog.json (tabelas geradas por
% docs/backlog/gera_backlog.py), docs/triagem/sintese/tematica/A4.md e B1 a B4.
% ---------------------------------------------------------------------------
\begingroup
\setlength{\emergencystretch}{3em}

\chapter{Requisitos e \textit{backlog} do protótipo}
\label{cap:requisitos}

Este capítulo junta as duas fontes de requisitos do buscador: o levantamento
com o especialista (Seção~\ref{sec:metodologia-requisitos}) e a síntese da
RSL~1 (Seção~\ref{sec:resultados-rsl1}). Em seguida, organiza os requisitos em
épicos e histórias de usuário e os prioriza para a entrega deste trabalho.

% ===========================================================================
\section{Requisitos acrescentados pela revisão}
\label{sec:requisitos-novos}

A síntese da RSL~1 consolidou os requisitos candidatos dos estudos em 27
requisitos e comparou cada um com o levantamento inicial
(Tabela~\ref{tab:rastreabilidade}). Dezoito já tinham um requisito
correspondente. Os outros nove não tinham. Seis deles viraram requisitos
novos (Tabela~\ref{tab:req-novos}). Os outros três (coleta na dark web,
classificação automática de conteúdo e suporte a vários idiomas) ficaram
fora desta entrega, como explica a Seção~\ref{sec:requisitos-priorizacao}.

% Gerado por docs/backlog/gera_backlog.py. Não editar à mão.
\begin{table}[htbp]
\centering
\caption{Requisitos novos propostos a partir da RSL 1}
\label{tab:req-novos}
\small
\begin{tabular}{lp{9.2cm}p{3cm}}
\hline
\textbf{ID} & \textbf{Requisito} & \textbf{Origem} \\
\hline
RF11 & Busca e filtros: buscar entidades e evidências do caso por nome, identificador e tipo. & RSL 1, A4 TD14 \\
RF12 & Alerta como pista: todo alerta exige confirmação ou descarte pelo investigador, com justificativa registrada na auditoria. & RSL 1, A4 TD18 e TA3 \\
RF13 & Extração assistida por LLM: entidades e campos extraídos de documentos, sempre com trecho e fonte citados e validação de esquema. & RSL 1, A4 TD13 \\
RF14 & Priorização de alertas: ordenar por severidade, com a política de erro (falso positivo ou falso negativo) declarada por tipologia. & RSL 1, A4 TD11 \\
RF15 & Análise temporal e geográfica: linha do tempo e mapa dos eventos do caso. & RSL 1, A4 TD15 \\
RNF12 & Resiliência da coleta: retentativa, controle de taxa, registro de falha por fonte e detecção de mudança de formato. & RSL 1, A4 TD4; A3 TA3 \\
\hline
\end{tabular}
\fonte{Elaborada pelo autor, a partir da síntese da RSL 1 (Seção~\ref{sec:resultados-rsl1}).}
\end{table}

\begin{table}[htbp]
\centering
\caption{Histórias do \textit{backlog} por prioridade MoSCoW}
\label{tab:moscow-resumo}
\small
\begin{tabular}{lrr}
\hline
\textbf{Prioridade} & \textbf{Histórias} & \textbf{Feitas no protótipo} \\
\hline
Must have & 11 & 11 \\
Should have & 13 & 5 \\
Could have & 24 & 8 \\
Won't have (nesta entrega) & 8 & 1 \\
\hline
\end{tabular}
\fonte{Elaborada pelo autor. MoSCoW de H40 a H57: decisão do autor, a validar pelo orientador.}
\end{table}

\begin{small}
\begin{longtable}{lp{7.9cm}p{1.3cm}cp{1.9cm}}
\caption{\textit{Backlog} do protótipo}\rotulounico{tab:backlog}\\
\hline
\textbf{ID} & \textbf{História} & \textbf{Req.} & \textbf{MoSCoW} & \textbf{Situação} \\ \hline\endhead
\multicolumn{5}{l}{\textbf{E1. Coleta de dados abertos}} \\
H01 & Como investigador, quero coletar empresas e sócios dos dados abertos do CNPJ e contratos do PNCP, para reunir num caso dados oficiais de fontes diferentes. & RF01, RNF10 & M & feito \\
H15 & Como investigador, quero saber, em cada coleta, quais fontes falharam e por quê, com retentativa automática, para não tirar conclusões de dados incompletos. & RNF12 & S & parcial \\
H19 & Como investigador, quero coletar também do Portal da Transparência e do TSE, para cruzar contratos com sanções e doações de campanha. & RF01 & C & feito \\
H25 & Como investigador, quero importar outros dados obtidos por ordem judicial (RIF do COAF, extratos telefônicos), separados do dado aberto, para analisá-los no mesmo caso. & RF10 & W & fora desta entrega \\
H27 & Como investigador, quero coletar de serviços da dark web, para investigar mercados ilícitos. & novo & W & fora desta entrega \\
H50 & Como investigador, quero coletar os bens declarados ao TSE pelos candidatos do caso em cada eleição, com a custódia de cada arquivo, para montar a linha do tempo patrimonial. & RF01, RF05 & S & aguarda orientador \\
H54 & Como investigador, quero importar dados bancários recebidos por ordem judicial no leiaute oficial do Banco Central, com validação de cada arquivo e conferência do período e das contas autorizados, para analisá-los no caso sem digitação. & RF10 & C & a fazer \\
\multicolumn{5}{l}{\textbf{E2. Entidades e vínculos}} \\
H02 & Como investigador, quero que a mesma empresa ou pessoa seja unificada entre as fontes, com sugestão (e não fusão) quando só o nome coincide, para não confundir homônimos. & RF02 & M & feito \\
H03 & Como investigador, quero ver o grafo de vínculos e o caminho entre duas entidades, para identificar sócios ocultos e cadeias societárias. & RF03 & M & feito \\
H41 & Como investigador, quero um alerta quando a mesma pessoa administrar muitas empresas abertas em pouco tempo, para identificar possível laranja profissional. & RF03 & S & a fazer \\
H42 & Como investigador, quero um alerta quando o controle de uma empresa do caso passar por várias pessoas jurídicas ou terminar em sócio estrangeiro, para identificar cadeias societárias em camadas. & RF03 & S & a fazer \\
\multicolumn{5}{l}{\textbf{E3. Alertas e análise assistida}} \\
H04 & Como investigador, quero alertas de tipologia (T1 a T3) que mostrem a regra, o motivo e as evidências com a fonte, para entender por que cada alerta foi gerado. & RF04, RNF05 & M & feito \\
H12 & Como investigador, quero marcar cada alerta como confirmado ou descartado, com justificativa registrada na auditoria, para que o alerta seja tratado como pista e a decisão fique com um humano. & RF12 & S & feito \\
H13 & Como investigador, quero ver os alertas ordenados por severidade, com a política de erro de cada tipologia declarada, para priorizar o que analisar primeiro. & RF14 & S & feito \\
H21 & Como investigador, quero ver os eventos do caso numa linha do tempo e, quando houver endereço, num mapa, para perceber padrões temporais e geográficos. & RF15 & C & a fazer \\
H22 & Como investigador, quero extrair entidades de documentos com apoio de LLM, sempre com o trecho e a fonte citados, para acelerar a leitura sem perder a verificabilidade. & RF13, RNF05 & C & parcial \\
H28 & Como investigador, quero classificar automaticamente textos e imagens por tipo de crime, para triar grandes volumes. & novo & W & fora desta entrega \\
H30 & Como investigador, quero analisar conteúdo em vários idiomas, com gírias e jargão criminal, para não perder informação em fontes estrangeiras. & novo & W & fora desta entrega \\
H31 & Como investigador, quero um alerta quando a empresa contratada não estiver com situação cadastral ativa no CNPJ, para identificar contratação irregular. & RF04 & S & feito \\
H32 & Como investigador, quero um alerta de possível fracionamento de compras (várias contratações do mesmo objeto, com o mesmo fornecedor, abaixo do limite de dispensa), para identificar fuga à licitação. & RF04 & C & feito \\
H33 & Como investigador, quero um alerta quando a empresa contratada estiver no cadastro de empresas sancionadas, para identificar contratação de empresa impedida. & RF04 & C & feito \\
H34 & Como investigador, quero um alerta quando um sócio da empresa contratada for servidor do órgão contratante, para identificar conflito de interesse. & RF04 & C & feito \\
H35 & Como investigador, quero um alerta de termos aditivos acima do limite legal, para identificar alteração indevida de contratos. & RF04 & C & feito \\
H36 & Como investigador, quero que o texto de editais seja analisado em busca de cláusulas restritivas, com o trecho citado, para identificar direcionamento da licitação. & RF13 & C & parcial \\
H39 & Como investigador, quero um alerta quando um doador de campanha (ou empresa de um sócio doador) for contratado pelo ente do candidato eleito, para identificar possível retribuição de doação. & RF04 & C & feito \\
H40 & Como investigador, quero um alerta quando muitas empresas estiverem registradas no mesmo endereço, mesmo sem contrato público, para identificar possível domicílio de fachada. & RF04, RF03 & S & a fazer \\
H43 & Como investigador, quero um alerta quando um MEI, microempresa ou empresa de pequeno porte receber em contratos públicos, no ano, mais que o limite de receita do seu regime, para identificar empresa de fachada. & RF04 & S & a fazer \\
H45 & Como investigador, quero um alerta quando uma empresa recém-aberta contratada tiver sócio que possivelmente é servidor, eleito ou PEP, para identificar desvio por empresa criada para contratar. & RF04 & C & a fazer \\
H46 & Como investigador, quero um alerta quando um eleito ou PEP for possivelmente sócio de empresa contratada pelo próprio ente durante o mandato, para identificar autocontratação. & RF04 & C & aguarda orientador \\
H47 & Como investigador, quero um alerta quando a contratada compartilhar sócio-administrador ou endereço com empresa sancionada, ou tiver sócio expulso da administração, para identificar empresa sucessora de sancionada. & RF04 & C & a fazer \\
H48 & Como investigador, quero um alerta de combinação quando a mesma entidade acumular vários alertas, com a lista de cada um e das evidências, para priorizar o que tem mais sinais juntos. & RF14, RF12 & C & a fazer \\
H49 & Como investigador, quero sinais fracos (contato compartilhado, sócio em faixa etária atípica, excesso de CNAEs) que só pesem no alerta de combinação, para não gerar alertas isolados de baixo valor. & RF04 & C & a fazer \\
H51 & Como investigador, quero ver a evolução dos bens declarados de um candidato entre eleições, e a renda pública conhecida quando houver, com alerta só quando a variação superar a renda, para identificar indício de enriquecimento ilícito. & RF05 & S & aguarda orientador \\
H52 & Como investigador, quero um alerta quando o candidato pagar, com recursos de campanha, empresa de atividade sem relação com campanha, recém-aberta ou com sinais de fachada, para identificar desvio de recursos de campanha. & RF04 & C & aguarda orientador \\
H53 & Como investigador, quero um alerta quando um servidor doar a campanha valor desproporcional à sua remuneração pública, para identificar possível doação por interposta pessoa. & RF04 & W & aguarda orientador \\
H56 & Como investigador, quero saber se o sócio de uma contratada recebe benefício social, para identificar possível laranja. & RF03 & W & fora desta entrega \\
H57 & Como investigador, quero um alerta quando uma entidade sem fins lucrativos que recebe convênio for recém-aberta, estiver no CEPIM ou tiver dirigente que possivelmente é eleito, PEP, servidor do concedente ou sócio de empresas do mesmo ramo, para identificar desvio por convênio. & RF04, RF03 & C & a fazer \\
\multicolumn{5}{l}{\textbf{E4. Prova e cadeia de custódia}} \\
H05 & Como investigador, quero que cada carga coletada seja guardada com hash SHA-256, carimbo de tempo UTC e origem, e que eu possa verificar a integridade depois, para que a prova se sustente em juízo. & RF07 & M & feito \\
H08 & Como investigador, quero exportar o caso (entidades, alertas, evidências e custódia) em JSON e Markdown, para anexar ao relatório de investigação. & RF09 & M & feito \\
H14 & Como perito, quero que a captura registre também a versão do conector e que eu possa recoletar a mesma fonte e comparar, para reconstituir como a prova foi obtida. & RF07 & S & feito \\
H24 & Como analista, quero exportar o caso num formato padronizado de troca (por exemplo, CASE/UCO), para levar os dados a outras ferramentas. & RNF09 & C & a fazer \\
H38 & Como analista, quero exportar o grafo do caso para i2 Analyst's Notebook e Maltego, para usar as ferramentas que o órgão já tem. & RNF09 & C & parcial \\
\multicolumn{5}{l}{\textbf{E5. Casos e busca}} \\
H07 & Como investigador, quero criar casos e fazer toda consulta dentro de um caso, com justificativa, para organizar a investigação e limitar a finalidade. & RF08 & M & feito \\
H09 & Como investigador, quero buscar entidades do caso por nome, CNPJ ou CPF parcial e filtrar por tipo, para encontrar rapidamente o que interessa. & RF11 & M & feito \\
H37 & Como investigador, quero anotar entidades e alertas e dividir tarefas do caso com a equipe, para trabalhar em conjunto na mesma investigação. & RF08 & C & feito \\
\multicolumn{5}{l}{\textbf{E6. Segurança, auditoria e conformidade}} \\
H06 & Como auditor, quero uma trilha só de acréscimo, encadeada por hash, de quem consultou o quê, quando e com qual justificativa, para detectar uso indevido. & RNF02 & M & feito \\
H10 & Como responsável pelo sistema, quero que todo acesso seja feito por usuário identificado e com perfil, para controlar quem usa a ferramenta. & RNF01 & M & feito \\
H11 & Como delegado responsável, quero que a ferramenta só faça coleta passiva de dados abertos, sem interagir com alvos, para não configurar infiltração virtual sem ordem judicial. & RNF04 & M & feito \\
H16 & Como desenvolvedor, quero análise estática de segurança e os testes como barreira no pipeline de integração, para que código gerado por IA não entre sem verificação. & RNF11 & S & parcial \\
H17 & Como delegado responsável, quero definir a finalidade e o prazo de retenção de cada caso, para aplicar a minimização dos dados coletados. & RNF04 & S & feito \\
H23 & Como responsável pelo sistema, quero MFA, acesso por caso e criptografia em repouso, para proteger dados de investigação. & RNF01 & C & parcial \\
H26 & Como investigador, quero coletar por infraestrutura não atribuível à polícia, para não alertar o alvo. & RNF03 & W & fora desta entrega \\
H55 & Como delegado responsável, quero que o dado sob ordem judicial fique numa base separada, com chave própria, marcado com o número da decisão, fora das exportações por padrão e com cada leitura auditada, para não misturá-lo ao dado aberto. & RF10, RNF02, RNF04 & C & a fazer \\
\multicolumn{5}{l}{\textbf{E7. Monitoramento e evolução}} \\
H18 & Como investigador, quero ser avisado quando mudar o quadro societário de uma empresa do caso ou surgir contrato novo, para acompanhar alvos sem refazer a coleta. & RF06 & C & feito \\
H29 & Como investigador, quero consultas de vínculo em segundos sobre grafos com milhões de nós, para trabalhar com a base completa do CNPJ. & RNF06 & W & feito \\
H44 & Como investigador, quero um alerta quando nenhum fundador continuar na empresa ou o quadro societário mudar pouco antes de um contrato, para identificar troca de laranjas. & RF03, RF06 & C & a fazer \\
\hline
\multicolumn{5}{p{14cm}}{\footnotesize Fonte: elaborada pelo autor. A evidência de cada história está em \texttt{docs/backlog/backlog.md}. MoSCoW de H40 a H57: decisão do autor, a validar pelo orientador.}\\
\end{longtable}
\end{small}


A revisão também refinou três requisitos existentes:
\begin{itemize}
	\item \textbf{RF07, cadeia de custódia:} a captura passa a registrar a
	versão do conector, e a mesma fonte pode ser recoletada e comparada, para
	reconstituir como a prova foi obtida.
	\item \textbf{RNF04, conformidade legal:} cada caso ganha finalidade e
	prazo de retenção definidos, o que aplica a minimização dos dados.
	\item \textbf{RNF05, explicabilidade:} toda saída de IA traz a fonte e,
	quando houver, o grau de confiança, e é tratada como pista, e não como
	conclusão.
\end{itemize}

A revisão não deu respaldo a parte do levantamento: RF04 (alertas de
tipologias de corrupção), RF05 (linha do tempo patrimonial), RF10
(integrações sob ordem judicial), as fontes brasileiras do RF01 e RNF11
(DevSecOps). Nenhum estudo incluído trata de corrupção em contratos
públicos nem do contexto brasileiro. Esses requisitos continuam no projeto,
justificados pelo levantamento com o especialista e pela legislação. O
RNF11 também tem apoio na RSL~2: o código gerado por IA não é seguro por
padrão, e as mitigações eficazes são ativas e em camadas
(Seção~\ref{sec:resultados-rsl2}).
% PENDENTE: o orientador (especialista e dono do produto) valida os
% requisitos novos RF11 a RF15 e RNF12 e os três refinamentos.

% ===========================================================================
\section{Comparação com ferramentas existentes e diferencial}
\label{sec:requisitos-mercado}

Para situar o protótipo, ele foi comparado com sete ferramentas, escolhidas
por decisão do autor em quatro grupos:
\begin{itemize}
	\item plataformas comerciais de análise de vínculos e OSINT voltadas a
	órgãos policiais e de inteligência: Maltego, i2 Analyst's Notebook e
	Palantir Gotham;
	\item uma ferramenta OSINT de código aberto: SpiderFoot;
	\item ferramentas do governo federal brasileiro para controle de
	contratações públicas: Alice, desenvolvida pela CGU e usada também pelo
	TCU, e o Observatório da Despesa Pública (ODP), da CGU;
	\item uma ferramenta de captura de evidências da internet para
	investigadores: OSIRT.
\end{itemize}
Os critérios de comparação vêm dos requisitos do buscador (Seção~\ref{sec:metodologia-requisitos}
e Seção~\ref{sec:requisitos-novos}). A fonte de cada célula é a documentação
oficial do fabricante ou do órgão, acessada em 13/11/2026. Essa documentação
é literatura cinzenta: descreve o produto como o fabricante o apresenta e não
passou por revisão por pares. Por isso, quando um recurso não aparece na
documentação, a tabela registra ``não encontrado'', e não que o produto não o
tenha. A comparação não incluiu testes práticos das ferramentas. Cada célula,
com o trecho literal e a página de origem, está em
\texttt{docs/mercado/concorrentes.json}.

% Gerado por docs/mercado/gera_mercado.py. Não editar à mão.
\begin{table}[htbp]
\centering
\caption[Comparação com ferramentas existentes: recursos]{Comparação do protótipo com ferramentas existentes: recursos}
\label{tab:mercado}
\footnotesize
\setlength{\tabcolsep}{2.5pt}
\begin{tabular}{lcccccccc}
\hline
\textbf{Ferramenta} & \textbf{Fontes BR} & \textbf{Entidades} & \textbf{Grafo} & \textbf{Tipologias} & \textbf{Evidência} & \textbf{Auditoria} & \textbf{IA} & \textbf{Português} \\
\hline
Maltego & -- & $\circ$ & $\bullet$ & -- & $\circ$ & $\circ$ & $\circ$ & -- \\
i2 ANB & -- & -- & $\bullet$ & -- & $\circ$ & -- & -- & -- \\
Gotham & -- & $\circ$ & $\bullet$ & -- & -- & $\bullet$ & $\circ$ & -- \\
SpiderFoot & -- & -- & $\circ$ & -- & -- & -- & -- & -- \\
Alice & $\bullet$ & -- & -- & $\bullet$ & -- & -- & $\bullet$ & $\bullet$ \\
ODP & $\circ$ & -- & $\circ$ & $\bullet$ & -- & -- & $\circ$ & $\bullet$ \\
OSIRT & -- & -- & -- & -- & $\bullet$ & $\circ$ & -- & $\circ$ \\
\textbf{Protótipo} & $\bullet$ & $\bullet$ & $\bullet$ & $\bullet$ & $\bullet$ & $\bullet$ & $\bullet$ & $\bullet$ \\
\hline
\end{tabular}
\fonte{Elaborada pelo autor, a partir da documentação oficial de cada produto, acessada em 13/11/2026 (lista em \texttt{docs/mercado/concorrentes.md}). $\bullet$ sim; $\circ$ parcial; -- não encontrado na documentação consultada, o que não significa que o produto não tenha o recurso.}
\end{table}

\begin{table}[htbp]
\centering
\caption[Comparação com ferramentas existentes: coleta e licença]{Comparação do protótipo com ferramentas existentes: coleta, implantação e licença}
\label{tab:mercado-modelo}
\small
\begin{tabular}{lp{3.4cm}p{3.4cm}p{4.2cm}}
\hline
\textbf{Ferramenta} & \textbf{Coleta} & \textbf{Implantação} & \textbf{Licença} \\
\hline
Maltego & parcial & ambos & comercial \\
i2 ANB & parcial & parcial & comercial \\
Gotham & -- & nuvem & comercial \\
SpiderFoot & ativa e passiva & local & aberto (MIT) \\
Alice & passiva & -- & uso interno \\
ODP & -- & -- & uso interno \\
OSIRT & passiva & local & comercial \\
Protótipo & só passiva & local (contêiner) & aberto (MIT) \\
\hline
\end{tabular}
\fonte{Elaborada pelo autor, a partir da documentação oficial de cada produto, acessada em 13/11/2026 (lista em \texttt{docs/mercado/concorrentes.md}). -- não encontrado na documentação consultada, o que não significa que o produto não tenha o recurso.}
\end{table}


As Tabelas~\ref{tab:mercado} e~\ref{tab:mercado-modelo} mostram três grupos de cobertura:
\begin{itemize}
	\item \textbf{Plataformas de análise de vínculos} (Maltego, i2 Analyst's
	Notebook e Gotham): cobrem o grafo de vínculos. A documentação consultada
	não mostra fontes de dados brasileiras integradas nem tipologias de fraude
	em contratações.
	\item \textbf{Ferramentas de controle da CGU} (Alice e ODP): cobrem o
	domínio. A Alice cruza os dados com ``mais de 40 bases de dados
	diferentes'' e aplica mais de 30 trilhas de auditoria (tipologias de
	risco). Mas são ferramentas de auditoria e controle interno: o manual
	informa que o acesso direto à Alice é restrito às unidades de auditoria
	interna e controle interno. A documentação consultada não descreve
	preservação da evidência com \textit{hash} nem trilha de auditoria das
	consultas, porque o produto do controle é o achado de auditoria, e não a
	prova de um inquérito.
	\item \textbf{Captura de evidências} (OSIRT): calcula o \textit{hash}
	SHA-512 e registra a data e a hora de cada captura. A documentação
	consultada não mostra grafo de vínculos nem tipologias.
\end{itemize}

A literatura mostra a mesma lacuna. Nenhum dos 57 estudos da RSL~1 trata de
fraude ou corrupção em contratações públicas. Só quatro juntam cadeia de
custódia e trilha de auditoria (RSL1-0466, RSL1-0816, RSL1-0854 e
RSL1-0863), e nenhum desses quatro tem grafo de vínculos
(Seção~\ref{sec:resultados-rsl1}).
% PENDENTE: a contagem por domínio usa o campo tipo_crime da extração, que
% ainda depende de R1 e R2.

\textbf{Diferencial.} Nenhuma ferramenta comparada nem nenhum estudo da
RSL~1 reúne, numa só ferramenta para a polícia judiciária:
\begin{enumerate}
	\item fontes de dados abertos brasileiras de contratações e de empresas,
	lidas no layout oficial;
	\item alertas de tipologias de fraude em contratações em que cada alerta
	traz a regra, o motivo e as evidências, e cada evidência aponta para a
	captura de origem;
	\item cadeia de custódia em cada captura, com \textit{hash} SHA-256,
	carimbo de tempo UTC e verificação da integridade a qualquer momento;
	\item trilha de auditoria só de acréscimo, encadeada por HMAC (\textit{hash} com chave), em
	que toda consulta é feita dentro de um caso e com justificativa;
	\item coleta só passiva por desenho, para não configurar infiltração
	virtual sem ordem judicial \cite[art. 10-A]{brasil2013lei12850};
	\item implantação local em contêineres, com código aberto sob a licença
	MIT, o que permite ao órgão auditar o código e manter os dados em
	infraestrutura própria.
\end{enumerate}
O diferencial está na combinação, voltada à produção de prova para o
inquérito. Isoladamente, cada recurso existe em alguma ferramenta.

\textbf{Limites da comparação.} O protótipo é muito menor que as ferramentas
comparadas: tem quatro fontes de dados abertos, além da leitura de editais, e dez tipologias, enquanto a Alice cruza mais
de 40 bases e aplica mais de 30 trilhas. A IA do protótipo é opcional e
restrita à extração de entidades de textos coletados: roda num servidor
local, cada sugestão cita o trecho literal de onde saiu e só entra no caso se
o analista aceitar (Capítulo~\ref{cap:desenvolvimento}). O protótipo não foi
avaliado com usuários policiais. A comparação mostra uma combinação que falta, e não que o
protótipo substitua essas ferramentas.

\textbf{Caminho para competir.} A comparação também mostra o que falta ao
protótipo. As histórias correspondentes estão no \textit{backlog}
(Tabela~\ref{tab:backlog}), em três horizontes:
\begin{enumerate}
	\item \textbf{Dentro deste trabalho} (histórias \textit{Should have}):
	confirmar ou descartar cada alerta (H12), ordenar os alertas por
	severidade (H13), registrar a versão do conector na captura (H14), relatar
	as falhas de coleta (H15), barrar código sem verificação no
	\textit{pipeline} (H16), definir a retenção por caso (H17) e alertar
	contratação de empresa sem situação cadastral ativa (H31), uma tipologia
	que a Alice também aplica. Além disso, avaliar o protótipo com usuários
	policiais, com questionário de usabilidade e tempo de tarefa comparado ao
	trabalho manual, como recomenda a síntese da QP5 da RSL~1. Sem essa
	avaliação, não há como afirmar que o protótipo se aproxima das
	ferramentas comparadas.
	\item \textbf{Aproximar-se das ferramentas de controle} (Alice e ODP), no
	domínio de contratações: mais fontes, como o Portal da Transparência e o
	TSE (H19); as tipologias do levantamento ainda não implementadas, como
	fracionamento de compras, empresa sancionada, sócio servidor e aditivos
	acima do limite legal (H32 a H35); monitoramento contínuo com aviso de
	novidades (H18), como os informes diários da Alice; e análise do texto
	dos editais com o trecho citado (H36).
	\item \textbf{Aproximar-se das plataformas de vínculos} (Maltego, i2 e
	Gotham): trabalho em equipe no caso (H37), exportação do grafo para i2 e
	Maltego (H38), MFA e controle de acesso por caso (H23), banco de grafos
	para escala (H29) e extração assistida por IA com a fonte citada (H22).
\end{enumerate}
O objetivo não é substituir essas plataformas. É ocupar o espaço que a
comparação mostrou vazio: uma ferramenta de código aberto, implantável pelo
próprio órgão, que produza alertas sobre contratações públicas já com a
evidência preservada para o inquérito.
% PENDENTE: o orientador valida os horizontes e as prioridades das histórias
% H31 a H38.

% ===========================================================================
\section{Épicos e histórias de usuário}
\label{sec:requisitos-historias}

Os requisitos foram escritos como histórias de usuário no formato ``Como
\textless papel\textgreater, quero \textless meio\textgreater, para
\textless fim\textgreater'', o mais difundido segundo
\citeonline[seção 1]{lucassen2016}. As histórias foram agrupadas em sete
épicos. Um épico é uma história grande, quebrada em histórias menores que
podem ser implementadas \cite[seção 2]{lucassen2016}. Cada história aponta o
requisito que atende e a evidência que a sustenta: o levantamento com o
especialista, a legislação ou a síntese da RSL~1 (para requisitos do
produto) ou da RSL~2 (para o processo de desenvolvimento). Os papéis são o
investigador, o delegado responsável, o auditor, o perito e o responsável
pelo sistema.

% ===========================================================================
\section{Priorização}
\label{sec:requisitos-priorizacao}

As histórias foram priorizadas pelo método MoSCoW, uma das técnicas de
priorização de requisitos mais usadas \cite{ais2026}. O método separa os
requisitos em quatro categorias: \textit{Must have}, \textit{Should have},
\textit{Could have} e \textit{Won't have} \cite[seção 1]{miranda2022}. Na
formulação usual, as três primeiras recebem 60\%, 20\% e 20\% do esforço de
desenvolvimento, e as de menor preferência só começam quando as de maior
preferência terminam \cite[seções 1 e 2]{miranda2022}. Para que essa reserva
proteja a entrega, cada categoria deve ter pelo menos cinco itens, e nenhum
deve consumir mais de 25\% do esforço da categoria \cite[seção 7]{miranda2022}.

A entrega considerada é a deste trabalho. Os critérios de classificação
foram decisão do autor:
\begin{itemize}
	\item \textbf{\textit{Must have}:} o que sustenta os dois eixos do
	levantamento, correlacionar dados abertos e produzir prova que se sustente
	em juízo, e o que a lei exige (coleta passiva).
	\item \textbf{\textit{Should have}:} o que a revisão pede e cabe no
	protótipo sem mudar a arquitetura, como confirmar ou descartar cada
	alerta e registrar a versão do conector na captura.
	\item \textbf{\textit{Could have}:} requisitos do levantamento ou da
	revisão que pedem nova fonte, novo componente ou mais esforço, como o
	monitoramento contínuo e a extração com LLM.
	\item \textbf{\textit{Won't have} (nesta entrega):} o que depende de dados
	sigilosos, de definição legal ainda aberta ou de outro domínio de crime,
	como as integrações sob ordem judicial, o anonimato operacional e a coleta
	na dark web.
\end{itemize}

O recorte do protótipo foi decidido antes da síntese da RSL~1, pelos dois
eixos do levantamento. Por isso, a priorização foi aplicada depois, sobre o
que já estava implementado. A divisão de esforço de 60, 20 e 20\% não foi
usada, porque o esforço de cada história não foi estimado antes. A
Tabela~\ref{tab:moscow-resumo} mostra quantas histórias há em cada
categoria, e a Tabela~\ref{tab:backlog} mostra o \textit{backlog} completo.
Das 39 histórias, 25 estão feitas, 6 estão em parte, 3 estão por fazer e 5
ficaram fora desta entrega (Tabela~\ref{tab:backlog}). As 11 histórias
\textit{Must have} estão feitas.
% PENDENTE: a priorização é proposta do autor. O orientador, que fez o
% levantamento inicial, valida as categorias como dono do produto.
% PENDENTE: quando as histórias Should have forem implementadas, atualizar a
% situação em docs/backlog/backlog.json e rodar docs/backlog/gera_backlog.py.

\endgroup
